\documentclass[10pt,letterpaper]{article}
\usepackage{technical-user-manual}
\setDocTitle{Chapter 19: Introduction to Assurance}
\setDocSubtitle{Book-aligned study guide and original verification exercise}
\setDocVersion{1.0}
\setDocStatus{Independent educational companion}
\setDocOwner{latex-docs / security}
\setDocDate{2026-11-02}
\setDocSummary{Independent guided study of Bishop Chapter 19, with an original falsifiable security exercise.}
\begin{document}
\makedoctocpage
\begin{notebox}[Source and authority boundary]
\textbf{BOOK:} Matt Bishop, \emph{Computer Security: Art and Science}, Second Edition (Pearson, 2019). Book-aligned sections are \emph{original summaries and reading prompts}, not copied textbook content.
\textbf{EXTENSION:} Any contemporary framework mappings, implementation tasks, scoring rules, and operational recommendations are independently authored educational supplements, not the book's positions. Nothing here proves compliance.
\end{notebox}
\section{Book-aligned reading scope}
\textbf{Chapter 19: Introduction to Assurance.} Read the corresponding chapter of Bishop (2019). The chapter's core study area is Specifications, trust, correctness arguments and evidence. This scope statement is a study organizer; no passage, diagram, proof, or exercise from the book has been reproduced here.
\section{Learning outcomes}
\begin{enumerate}
\item Explain the threat, security objective, and trust assumption motivating introduction to assurance.
\item Differentiate the formal policy or security property from an implementation mechanism and its operational evidence.
\item Evaluate a counterexample; explain whether it disproves a claim, violates an assumption, or lies outside the chosen model.
\item Write a reproducible negative test and identify the limit of the evidence it provides.
\end{enumerate}
\section{Security assertion and model boundary}
\begin{infobox}[Original analytical proposition]
A requirement needs a concrete test or proof obligation before assurance can be assessed
\end{infobox}
Construct a model with subjects, resources, relevant state, permitted transitions, and an attacker-observable outcome. State what is explicitly excluded from the model (for example, an operating-system compromise, physical access, or uncontrolled side channels). Do not treat a successful finite test as a general mathematical proof.
\section{Guided reading questions}
\begin{enumerate}
\item Which entities and operations must be represented before the chapter's central security question can be asked precisely?
\item What is the difference between a policy violation, a reachable exploit, and an audit finding?
\item Which assumptions would fail if a user or process were malicious rather than merely mistaken?
\item What counterexample would invalidate the proposed design, and how could it be observed in a laboratory?
\item Where does the chapter's formal or conceptual treatment stop short of a deployable security control?
\end{enumerate}
\section{Independent application exercise}
\textbf{Scenario:} Create a claim-argument-evidence tree for a privileged workflow.
\begin{enumerate}
\item Draw a state/actor/data-flow model and label trust boundaries.
\item Write a falsifiable claim and at least three test cases: a permitted operation, a prohibited operation, and a boundary or failure case.
\item Capture immutable source revision, test environment, exact commands, expected outcome and actual outcome.
\item Record false positives, false negatives, unresolved risks and an independently reviewable conclusion.
\end{enumerate}
\textbf{Assessment artifact:} Evidence list with version, scope, owner and gaps.
\section{Contemporary engineering bridge --- independent supplement}
Assurance cases and NIST SP 800-53 is a useful adjacent reference, not terminology or a requirement attributed to Bishop. Consult authoritative, dated source publications before citing exact framework identifiers. Map \emph{security claims} to evidence; a framework tag alone does not implement a control.
\section{Self-assessment and completion gate}
\begin{enumerate}
\item Explain the policy in a single unambiguous sentence without relying on product names.
\item Present a concrete counterexample or state why one is excluded by a specified assumption.
\item Demonstrate at least one passing and one intentionally rejected case in a non-production environment.
\item Submit a reviewer-ready evidence record; mark unsupported generalizations as limitations.
\end{enumerate}
\textbf{Pass criterion:} The claim, model boundary, negative result, evidence provenance and residual limitation are separately documented. Suggested competency rating: 0 = missing; 1 = described; 2 = demonstrated; 3 = independently verified.
\section{Further study and cross-references}
Return to Bishop's chapter bibliography for foundational scholarship. For implementation guidance, consult the repository's existing \texttt{src/security/security-engineering/}, \texttt{src/security/application-security/sast/}, and language-specific secure coding collections rather than creating parallel versions.

\end{document}
